%Version 3 December 2023
% See section 11 of the User Manual for version history
%
%%%%%%%%%%%%%%%%%%%%%%%%%%%%%%%%%%%%%%%%%%%%%%%%%%%%%%%%%%%%%%%%%%%%%%
%%                                                                 %%
%% Please do not use \input{...} to include other tex files.       %%
%% Submit your LaTeX manuscript as one .tex document.              %%
%%                                                                 %%
%% All additional figures and files should be attached             %%
%% separately and not embedded in the \TeX\ document itself.       %%
%%                                                                 %%
%%%%%%%%%%%%%%%%%%%%%%%%%%%%%%%%%%%%%%%%%%%%%%%%%%%%%%%%%%%%%%%%%%%%%

%%\documentclass[referee,sn-basic]{sn-jnl}% referee option is meant for double line spacing

%%=======================================================%%
%% to print line numbers in the margin use lineno option %%
%%=======================================================%%

%%\documentclass[lineno,sn-basic]{sn-jnl}% Basic Springer Nature Reference Style/Chemistry Reference Style

%%======================================================%%
%% to compile with pdflatex/xelatex use pdflatex option %%
%%======================================================%%

%%\documentclass[pdflatex,sn-basic]{sn-jnl}% Basic Springer Nature Reference Style/Chemistry Reference Style


%%Note: the following reference styles support Namedate and Numbered referencing. By default the style follows the most common style. To switch between the options you can add or remove Numbered in the optional parenthesis. 
%%The option is available for: sn-basic.bst, sn-vancouver.bst, sn-chicago.bst%  
 
%%\documentclass[pdflatex,sn-nature]{sn-jnl}% Style for submissions to Nature Portfolio journals
%%\documentclass[pdflatex,sn-basic]{sn-jnl}% Basic Springer Nature Reference Style/Chemistry Reference Style


\documentclass[pdflatex,sn-mathphys-num]{sn-jnl}% Math and Physical Sciences Numbered Reference Style 

% TODO: uncomment the line above, the line below is for double-spacing only
% \documentclass[referee,pdflatex,sn-mathphys-num]{sn-jnl}


%%\documentclass[pdflatex,sn-mathphys-ay]{sn-jnl}% Math and Physical Sciences Author Year Reference Style
%%\documentclass[pdflatex,sn-aps]{sn-jnl}% American Physical Society (APS) Reference Style
%%\documentclass[pdflatex,sn-vancouver,Numbered]{sn-jnl}% Vancouver Reference Style
%%\documentclass[pdflatex,sn-apa]{sn-jnl}% APA Reference Style 
%%\documentclass[pdflatex,sn-chicago]{sn-jnl}% Chicago-based Humanities Reference Style

%%%% Standard Packages
%%<additional latex packages if required can be included here>

\usepackage{graphicx}%
\usepackage{multirow}%
\usepackage{amsmath,amssymb,amsfonts}%
\usepackage{amsthm}%
\usepackage{mathrsfs}%
\usepackage[title]{appendix}%
\usepackage{xcolor}%
\usepackage{textcomp}%
\usepackage{manyfoot}%
\usepackage{booktabs}%
\usepackage{algorithm}%
\usepackage{algorithmicx}%
\usepackage{algpseudocode}%
\usepackage{listings}%
\usepackage{hyperref}%
\usepackage{enumitem}
%%%%

%%%%%=============================================================================%%%%
%%%%  Remarks: This template is provided to aid authors with the preparation
%%%%  of original research articles intended for submission to journals published 
%%%%  by Springer Nature. The guidance has been prepared in partnership with 
%%%%  production teams to conform to Springer Nature technical requirements. 
%%%%  Editorial and presentation requirements differ among journal portfolios and 
%%%%  research disciplines. You may find sections in this template are irrelevant 
%%%%  to your work and are empowered to omit any such section if allowed by the 
%%%%  journal you intend to submit to. The submission guidelines and policies 
%%%%  of the journal take precedence. A detailed User Manual is available in the 
%%%%  template package for technical guidance.
%%%%%=============================================================================%%%%

%% as per the requirement new theorem styles can be included as shown below
\theoremstyle{thmstyleone}%
\newtheorem{theorem}{Theorem}%  meant for continuous numbers
%%\newtheorem{theorem}{Theorem}[section]% meant for sectionwise numbers
%% optional argument [theorem] produces theorem numbering sequence instead of independent numbers for Proposition
\newtheorem{proposition}[theorem]{Proposition}% 
%%\newtheorem{proposition}{Proposition}% to get separate numbers for theorem and proposition etc.

\theoremstyle{thmstyletwo}%
\newtheorem{example}{Example}%
\newtheorem{remark}{Remark}%

\theoremstyle{thmstylethree}%
\newtheorem{definition}{Definition}%

\raggedbottom
%%\unnumbered% uncomment this for unnumbered level heads

\begin{document}

\title[Article Title]{Noise-Robust QSAR: Evaluating Probabilistic Methods Across Molecular Representations}
%%=============================================================%%
%% GivenName	-> \fnm{Joergen W.}
%% Particle	-> \spfx{van der} -> surname prefix
%% FamilyName	-> \sur{Ploeg}
%% Suffix	-> \sfx{IV}
%% \author*[1,2]{\fnm{Joergen W.} \spfx{van der} \sur{Ploeg} 
%%  \sfx{IV}}\email{iauthor@gmail.com}
%%=============================================================%%

\author[1]{\fnm{Adelaide} \sur{Punt}}\email{adelaidepunt@gmail.com}

\author[2]{\fnm{Thierry} \sur{Hanser}}

\author[2]{\fnm{Stephane} \sur{Werner}}

\author*[1]{\fnm{Garrett} \sur{Morris}}\email{garrett.morris@stats.ox.ac.uk}

\affil*[1]{\orgdiv{Department of Statistics}, \orgname{University of Oxford}, \orgaddress{\street{24-29 St Giles'}, \city{Oxford}, \postcode{OX1 3LB}, \country{UK}}}

\affil[2]{\orgdiv{Lhasa Limited}, \orgname{Granary Wharf House}, \orgaddress{\street{2 Canal Wharf}, \city{Leeds}, \postcode{LS11 5PS}, \country{UK}}}

%%==================================%%
%% Sample for unstructured abstract %%
%%==================================%%
\abstract{Predictive models for chemical bioactivity are limited by experimental noise and biological variability, yet most machine learning approaches treat noisy labels as ground truth. We systematically evaluated how molecular representation and model architecture affect both prediction accuracy and noise robustness, comparing 14 model architectures across 6 representations and 6 noise injection strategies on the QM9 dataset and three external ADME datasets. Using ANOVA variance decomposition, we demonstrate a role inversion: representation and model contribute comparably to prediction performance, but under label noise, model architecture becomes the dominant factor. Three distinct robustness mechanisms emerge: inherent robustness through inductive biases (SVM, full Bayesian neural networks), representation-mediated robustness that depends on input features (RF, MLP), and ensemble-based robustness from dataset-specific structure. Among probabilistic models, those whose loss functions provide an explicit channel for noise absorption---routing noise into uncertainty estimates rather than distorting predictions---best detect and resist noise. Model rankings are highly concordant across noise types but do not transfer to external datasets, with transferability depending on the robustness mechanism rather than the ranking itself. We introduce NoiseInject, an open-source framework for benchmarking noise robustness across arbitrary datasets.}


%%================================%%
%% Sample for structured abstract %%
%%================================%%

% \abstract{\textbf{Purpose:} The abstract serves both as a general introduction to the topic and as a brief, non-technical summary of the main results and their implications. The abstract must not include subheadings (unless expressly permitted in the journal's Instructions to Authors), equations or citations. As a guide the abstract should not exceed 200 words. Most journals do not set a hard limit however authors are advised to check the author instructions for the journal they are submitting to.
% 
% \textbf{Methods:} The abstract serves both as a general introduction to the topic and as a brief, non-technical summary of the main results and their implications. The abstract must not include subheadings (unless expressly permitted in the journal's Instructions to Authors), equations or citations. As a guide the abstract should not exceed 200 words. Most journals do not set a hard limit however authors are advised to check the author instructions for the journal they are submitting to.
% 
% \textbf{Results:} The abstract serves both as a general introduction to the topic and as a brief, non-technical summary of the main results and their implications. The abstract must not include subheadings (unless expressly permitted in the journal's Instructions to Authors), equations or citations. As a guide the abstract should not exceed 200 words. Most journals do not set a hard limit however authors are advised to check the author instructions for the journal they are submitting to.
% 
% \textbf{Conclusion:} The abstract serves both as a general introduction to the topic and as a brief, non-technical summary of the main results and their implications. The abstract must not include subheadings (unless expressly permitted in the journal's Instructions to Authors), equations or citations. As a guide the abstract should not exceed 200 words. Most journals do not set a hard limit however authors are advised to check the author instructions for the journal they are submitting to.}

\keywords{QSAR, noise robustness, uncertainty quantification, label noise, molecular representations, Bayesian neural networks, molecular property prediction}

%%\pacs[JEL Classification]{D8, H51}

%%\pacs[MSC Classification]{35A01, 65L10, 65L12, 65L20, 65L70}

\maketitle

\noindent\textbf{Scientific contribution.} This work introduces ANOVA variance decomposition as a tool for attributing noise robustness to model architecture versus molecular representation in QSAR---revealing that these two factors trade dominance between prediction performance and noise robustness, a pattern we term the role inversion. We identify three distinct mechanisms by which models resist label noise and show that only models whose loss functions provide explicit noise-absorption channels produce robustness that transfers across datasets. The accompanying NoiseInject framework enables systematic noise robustness benchmarking on arbitrary molecular property datasets.


\section*{Abbreviations}

\begin{description}[style=unboxed,leftmargin=0cm]
    \item[ADME] Absorption, distribution, metabolism, and excretion
    \item[ANOVA] Analysis of variance
    \item[BNN] Bayesian neural network
    \item[DFT] Density functional theory
    \item[DNN] Deep neural network
    \item[ECFP] Extended-connectivity fingerprint
    \item[ECE] Expected calibration error
    \item[ELBO] Evidence lower bound
    \item[GNN] Graph neural network
    \item[GP] Gaussian process
    \item[HOMO] Highest occupied molecular orbital
    \item[ICC] Intraclass correlation coefficient
    \item[LightGBM] Light Gradient Boosting Machine
    \item[LNL] Learning with noisy labels
    \item[LUMO] Lowest unoccupied molecular orbital
    \item[MAE] Mean absolute error
    \item[MHG-GNN] Molecular Hypergraph Grammar Graph Neural Network
    \item[ML] Machine learning
    \item[MLP] Multi-layer perceptron
    \item[NDS] Noise degradation slope
    \item[NGBoost] Natural gradient boosting
    \item[NN] Neural network
    \item[PDV] Physicochemical descriptor vector
    \item[PLS] Partial least squares
    \item[QRF] Quantile regression forest
    \item[QSAR] Quantitative structure--activity relationship
    \item[QSPR] Quantitative structure--property relationship
    \item[RF] Random forest
    \item[RMSE] Root mean squared error
    \item[SMILES] Simplified molecular input line entry system
    \item[SNS] Sort \& Slice
    \item[SVM] Support vector machine
    \item[VBLL] Variational Bayesian Last Layers
    \item[XGBoost] eXtreme Gradient Boosting
\end{description}

\clearpage

\section{Introduction}\label{sec1}

New avenues in the drug discovery pipeline have opened thanks to advances in large data collection, storage, and analysis. Therapeutic targets are identified, potential drug candidates are screened and optimized, and the most promising of those go on to pre-clinical and clinical trials \citep{Lill2007}. There are various methods by which potential drug candidates can be screened. Virtual screening, which allows for the rapid and large-scale analysis of libraries of data, is a particularly useful approach used to speed up the drug discovery process \citep{Neves2018}. One standard screening method is the prediction of chemical bioactivities and physical properties of such candidates. Quantitative structure-activity relationship (QSAR) models can rapidly make these predictions, act as a filter for compounds with undesirable properties, and thus reduce the time and cost associated with experimental screening \citep{Sabando2021}. Specifically, QSAR models predict the relationship between molecular features encoded within a particular molecular representation and a target property \citep{Sabando2021}. A QSAR model typically consists of a molecular representation paired with a machine learning (ML) model. An ideal molecular representation will allow a model to learn only the relevant information about each molecule with respect to the target property. Molecular featurization refers to the process of converting a molecule into a computational representation, or an embedding \citep{Sabando2021}. 

QSAR models are sensitive to the quality of the input data on which they are trained, in the form of experimental noise from measurement errors, variability under experimental conditions, provenance of the data and intrinsic biological and chemical variability. \cite{Kolmar2021} show that the variance associated with the experimental data can contribute more to the prediction error than the model's error. Even measures of model performance may be flawed if test and validation sets are laden with noise. In the context of activity prediction, the inherent variability within experimental setups and assays, batch effect, measurement error, and biological variability introduce additional layers of complexity \citep{Kolmar2021, Fourches2010}. All of this increases the likelihood of encountering noise in the data \citep{Fang2022}. Even the representations themselves may be misleading, as they encode structural information and may not account for all the potential conformations a molecule can adopt, primarily when binding to a particular target. 

Numerous studies have explored learning with noisy labels (LNL) in the broader context of ML and specifically within QSAR models. \cite{Song2022} provides a comprehensive review of deep LNL, identifying five main approaches: modifying architectures to better handle noise in training, employing regularization to reduce overfitting to noise, adjusting loss functions to be less sensitive to outliers or label flips, post-processing to adjust the weights for noisy data, and sample selection to select presumably clean data to pre-train. Although some QSAR models can overcome noise, predictive performance depends on the algorithm. \cite{Cortes2015} examined 12 different algorithms across 12 datasets and 10 noise levels, finding that no single algorithm was overall best at noise; certain algorithms perform better or worse at different levels of noise. This is a recurring theme in LNL research: There is no universal solution to handling noise in cheminformatics data. This is consistent with the No Free Lunch Theorem, which states that there is no single learning algorithm that can outperform all others in all data distributions \citep{Wolpert1997}. This suggests that there is no single best model or molecular representation for noise robustness; it depends on the underlying data. Certain techniques can be more effective with different types of noise. For example, small-loss selection works particularly well for symmetric noise \citep{Song2022}. Some techniques that may be robust to experimental noise may not be fit for the particular QSAR task at hand. Combinations of LNL approaches have yielded promising results, specifically sample selection combined with semi-supervised learning on rejected data, presumably noisy samples. Many LNL techniques were developed for image classification tasks. The data properties of molecular datasets differ significantly, so adjustments are needed to adapt some of these techniques \citep{Song2022}.

One key issue with QSAR models is the standard ML assumption that labels represent true values, as most chemical data sets lack sufficient replications to provide a robust statistical representation of the underlying distributions \citep{Kolmar2021}. However, probabilistic methods such as Gaussian processes (GPs) and Bayesian neural networks (BNNs) do not make this assumption \citep{Kolmar2021}. Unlike standard deterministic ML methods, Bayesian methods model distributions over predictions using Bayes' rule, allowing for the quantification of uncertainty within parameters and labels.

One ML method that has had particular success in molecular property prediction is Gaussian Processes (GPs) \citep{Obrezanova2007, gauche}. GPs have been known to match or exceed the predictive performance of traditional QSAR models such as support vector machines (SVMs) or partial least squares (PLS), while also providing insight into the reliability of those predictions \citep{Obrezanova2007}.

Individual decision trees tend to have high variance and often overfit, particularly at large depths \citep{Breiman1984}. Small modifications made to the training data can result in large differences in the structure of the tree, resulting in unstable predictions \citep{Breiman2001}. Ensemble methods are used to counteract this instability, aggregating predictions from many trees. Bootstrapping ensures that each tree is trained on a slightly different subset of data, which effectively reduces the influence of any individual sample \citep{Breiman2001}. This is especially beneficial with noisy samples.

NNs are a popular choice in QSAR research, though they are often outperformed by more traditional methods \citep{Baskin2008, Koutsoukas2017}. NNs are often deterministic; however, they can be transformed into probabilistic models through several approaches. Bayesian approaches also act as forms of regularization \citep{Burden2009}, as dropouts help to prevent overfitting.

Although numerous studies have explored the effects of noise on ML algorithms for molecular property prediction, comparisons between different molecular representations remain limited. Kolmar and Unsleber \cite{Kolmar2021} examined the predictive limits of QSAR in the presence of noise, showing that by learning the underlying trends rather than fitting to noise, the models can achieve predictions more accurate than those of their training data. They assumed experimental noise could be modeled as Gaussian-distributed, using the Central Limit Theorem to argue that aggregated experimental measurements from multiple independent sources of error exhibit a Gaussian distribution. However, Heid et al.\ \cite{Heid2023} used mean-variance estimation and bias-variance decomposition to show that noise can be tied to specific modalities, for example structure dependence. As an example, nitrogen-containing functional groups introduce measurement deviations, implying the presence of systematic noise. The choice of molecular representation also impacts the predictive performance of QSAR models. Deng et al.\ \cite{Deng2023} conducted a large-scale evaluation of 62,820 models and found that Extended Connectivity Fingerprint (ECFP) fingerprints frequently outperform learned graph neural network (GNN) representations, with dataset size proving more important than representation choice.

This research addresses three main questions concerning the robustness of molecular machine learning models under label noise. First, we investigate the contributions of molecular representation and model architecture to both overall prediction performance and, more specifically, noise robustness. Second, we compare how noise-robust probabilistic models are to their deterministic counterparts, and whether uncertainty estimates correlate with prediction error or label noise under noisy conditions. Third, we assess the generalizability of noise-robustness patterns across different noise-injection mechanisms and molecular properties. Together, these questions aim to identify what, if anything, makes a QSAR model robust to noise. Finally, we have released NoiseInject, an open-source Python package that acts as a noise benchmarking tool. It allows the user to add noise in different scenarios to an arbitrary dataset and to compute metrics to analyze its effects.

\section{Methods}\label{sec2}

\subsection{Dataset}
The majority of experiments were conducted on the QM9 molecular property dataset \citep{Ramakrishnan2014}, which contains approximately 130,000 small organic molecules with pre-computed quantum-mechanical properties obtained from density functional theory (DFT) calculations at the B3LYP/6-31G(2df,p) level of theory \citep{Ramakrishnan2014}. Each machine learning experiment was replicated 10 times with a distinct random seed and a subset of $N = 10{,}000$ molecules. Experiments that involved tracking uncertainty values were only run once due to computational limitations. For each experimental replicate, a new scaffold-based train/validation/test split (80/10/10) implemented by DeepChem \citep{deepchem} was generated from this fixed subset. Scaffold splits were used to prevent structural overlap between sets. When switching from random splitting to a more challenging scenarios such as scaffold splitting, QSAR models typically see performance drops \citep{Heid2023}.

We selected HOMO-LUMO energy gap as the primary prediction target in QM9, as it captures electronic characteristics relevant to molecular reactivity and stability \citep{Islam2019, Hllermeier2021}. The HOMO-LUMO gap is a well-established target for quantum chemistry benchmarks as it is defined as the difference between the highest occupied and lowest unoccupied molecular orbital energies, directly relating to a molecule's electronic excitability, charge transfer capability, and chemical stability \citep{Fediai2023}. We also assumed that thanks to the consistent calculation method, and notwithstanding approximations in the level of theory used, the QM data were "free" of noise.

To assess whether noise robustness results generalize beyond the HOMO-LUMO gap, we evaluated model configurations across three additional molecular property datasets with experimentally measured endpoints. From the OpenADMET benchmark suite, we used LogD (lipophilicity, $N = 5{,}039$) and Caco-2 efflux permeability ($N = 2{,}161$), both of which are key ADME properties in drug discovery. From ChEMBL, we used hERG-Ki ($N = 1{,}482$), measuring binding affinity to the cardiac ion channel hERG (Kv11.1) as a continuous regression target (pKi values) obtained by filtering for binding assays (type `B'), deduplicating by median pChEMBL value per compound, and removing compounds with standard deviation $> 1.0$ log unit. All three datasets were evaluated using all six noise injection strategies and the same noise levels ($\sigma \in \{0, 0.1, \ldots, 1.0\}$) as the QM9 experiments, enabling direct comparison of robustness rankings across prediction targets.

Prior to modeling, all molecular structures were sanitized and canonicalized using \texttt{RDKit} to remove invalid valence states and standardize atom and bond typing \citep{rdkit}. Target values were mean-centered and normalized to zero mean and unit variance. Data loading and pre-processing were performed with a global random seed to ensure reproducibility.

\subsection{Molecular Representations}

In this study, we evaluated a variety of molecular representations spanning path-based fingerprints, circular fingerprints, physicochemical descriptors, string encodings, and pretrained embeddings. We used three binary fingerprints: an RDKit topological fingerprint, a Morgan circular fingerprint, and the \emph{Sort \& Slice} (SNS) fingerprint \citep{Dablander2024}. The topological fingerprint enumerates all linear and branched paths through the molecular graph up to length 7 bonds, hashing them into a $d=2048$-bit vector using \texttt{RDKit} \citep{rdkit}. Each bit encodes the presence of one or more path-based substructures. The Morgan fingerprint uses circular substructures with radius $r=2$ (equivalent to ECFP4) hashed into a $d=2048$-bit vector, capturing the local chemical environment around each atom rather than linear paths. One downside of hashed fingerprints is bit collisions, in which chemically distinct substructures map to the same index, resulting effectively in representation-level noise. To address this issue, the SNS fingerprint uses a collision-free alternative based on the same Morgan circular substructures with radius $r=2$ \citep{Dablander2024}. Substructures are sorted by prevalence in the training set, and the top $L$ ($L=1024$) are sliced to create a binary vector, providing the ECFP-based circular fingerprint in our comparison. We also used 200-dimensional physicochemical descriptor vectors (PDVs) computed from RDKit molecular descriptors, and one-hot-encoded (OHE) SMILES strings \citep{Bertoni2020}, generated with \texttt{RDKit} \citep{rdkit}. We included mol2vec embeddings \citep{jaeger2018mol2vec}, which apply a Word2Vec-inspired approach to molecular substructures. Molecules are decomposed into Morgan substructures at multiple radii, treated as ``words'' in a molecular ``sentence,'' and a pre-trained Word2Vec model produces 300-dimensional embeddings for each substructure. The molecular embedding is the mean of its constituent substructure vectors, yielding a fixed-length, continuous representation that captures substructural context without requiring task-specific training. Finally, we included graph embeddings generated by Molecular Hypergraph Grammar GNNs (MHG-GNN) \citep{kishimoto2023}. MHG-GNNs are a GIN-based autoencoder pretrained on 1.34 million PubChem molecules using $\beta$-VAE loss, producing 1024-dimensional embeddings through iterative message passing. MHG-GNN was originally developed for materials science, although it has demonstrated strong performance on property prediction tasks for polymers, photoresists, and chromophores \citep{kishimoto2023}. 

All representations are precomputed through a Rust–Python pipeline. Rust performs various heavy-computation tasks involved in feature extraction, normalization, and serialization, and outputs the results into memory-mapped binaries (\texttt{.mmap}) for zero-copy reads. Model training, testing, and validation are then done in Python, reading in from the memory-mapped files.

\subsection{Models}
In this work, we used a wide range of model architectures to understand both deterministic and probabilistic learning behaviors across molecular representations and levels of noise. One key issue with QSAR models is the standard ML assumption that labels represent true values, as most chemical data sets lack sufficient replications to provide a robust statistical representation of the underlying distributions \citep{Kolmar2021}. A single measurement, or even multiple measurements for a given label, will not guarantee that this label accurately reflects the population mean \citep{Kolmar2021}. Most ML models rely on the assumption that the training data is, in fact, an accurate reflection of the population. Therefore, they treat these potentially noisy labels as discrete quantities rather than distributions, which opens the door to overfitting to noise and giving unduly optimistic assessments of the models. However, probabilistic methods such as Gaussian processes (GPs) and Bayesian neural networks (BNNs) do not make this assumption \citep{Kolmar2021}.

Across all models, hyperparameter tuning was performed using Optuna \citep{akiba2019optuna} Bayesian optimization (Tree-structured Parzen Estimator). The tuned hyperparameters were compared against conventional default values for each architecture, and the configuration with better validation performance was retained as the fixed parameter set for all subsequent experiments (Supplementary Table~\ref{tab:supp_s1_hyperparameters}).

The uncertainty values derived from these non-deterministic methods can be further decomposed into epistemic and aleatoric components. Epistemic uncertainty accounts for uncertainty within the model and can be explained away with sufficient data. Aleatoric uncertainty refers to the noise inherent in the data itself, meaning more data will not improve aleatoric uncertainty \citep{Hllermeier2021, kendall2017}.

We included the following variants of decision trees: Random Forest (RF), Quantile Regression Forest (QRF), eXtreme Gradient Boosting (XGBoost), Light Gradient Boosting Machine (LightGBM), and Natural Gradient Boosting (NGBoost). RFs are one of the most common choices for QSAR modeling thanks to their robustness and interpretability \citep{Svetnik2003}. However, variance across tree predictions can be used as an uncertainty estimate, and various QSAR studies have linked this ensemble variance to the prediction error; it is not as robust a quantification of uncertainty as found in probabilistic modeling  \citep{Svetnik2003}.  QRFs extend these RF predictions to full distributions \citep{Meinshausen2006}. QRFs keep the distribution of training labels in each leaf and, when they make a prediction, they use all the training instances in the leaves that the new instance falls into, across all trees, to compute quantiles from that set of values. This results in a non-parametric estimate of the conditional distribution of the output \citep{Meinshausen2006}. It inherently accounts for heteroscedasticity, such that if the target variable has more spread in its domain, the quantiles will reflect that. QRF has also been applied in QSAR, though not as widely as RF \citep{Venkatraman2021}. We included XGBoost, a standard among molecular property prediction tasks \citep{Mustapha2016, Tian2022}. We also included LightGBM \citep{ke2017lightgbm}, a gradient boosting framework that uses histogram-based split finding and leaf-wise tree growth for efficient training. LightGBM employs Gradient-based One-Side Sampling (GOSS) and Exclusive Feature Bundling (EFB) to handle high-dimensional features efficiently \citep{ke2017lightgbm}. Alongside XGBoost and LightGBM, we used Natural Gradient Boosting (NGBoost) \citep{Duan2020} that extends traditional gradient-boosted methods by outputting predictive distributions rather than point estimates. NGBoost treats the parameters of a chosen parametric distribution as regression targets and learns them via boosting with a natural gradient update rule \citep{Duan2020}.

We also included Support Vector Machines (SVMs) with radial basis function (RBF) kernels \citep{Vapnik1995}. SVMs are a well-established baseline in QSAR modeling, mapping inputs into a high-dimensional feature space where a maximum-margin hyperplane separates predictions \citep{Svetnik2003}. Although SVMs produce only point predictions without native uncertainty quantification, they serve as a deterministic reference for comparison with probabilistic approaches.

One ML method that has had particular success in molecular property prediction are Gaussian Processes (GPs) \citep{Obrezanova2007, gauche}. A GP is a non-parametric model that outputs a Gaussian predictive distribution over every data point, providing an uncertainty estimate for each prediction \citep{Obrezanova2007}. GPs have been known to match or exceed the predictive performance of traditional QSAR models such as support vector machines (SVMs) or partial least squares (PLS), while also providing insight into the reliability of those predictions \citep{Obrezanova2007}. The GP defines a distribution over functions, whose structure is determined by the kernel. These kernels can be modified to fit specific domains or functionalities. Typically, a Tanimoto kernel is used when working with molecular fingerprints \citep{Ralaivola2005, moss2020, gauche}. One of the major limitations to using GPs in practice is computational scaling. Inference has $O(N^3)$ time complexity and $O(N^2)$ memory, where $N$ is the number of training points \citep{Rasmussen2005}. This can be mitigated with sparse GP approximations, but these introduce approximation error and can be more difficult to tune \citep{quinonero2005}. Ensembles of trees or BNNs are typically more efficient for larger data sizes \citep{lakshminarayanan2017}.

We also employed GP models using the \texttt{Gauche} framework \citep{gauche}, which is optimized for cheminformatics tasks. Standardized kernels for GPs are not optimized for the chemical landscape. \texttt{Gauche} defines such kernels for strings, fingerprints, and graphs that operate on a range of widely used molecular representations \citep{gauche}. In this study, we used the Tanimoto kernel, which is defined for binary vectors $\mathbf{x}, \mathbf{x}' \in \{0, 1\}^d$ for $d \geq 1$ as:

$$
k_{\text{Tanimoto}}(\mathbf{x}, \mathbf{x}') := \sigma_f^2 \cdot \frac{\langle \mathbf{x}, \mathbf{x}' \rangle}{\|\mathbf{x}\|^2 + \|\mathbf{x}'\|^2 - \langle \mathbf{x}, \mathbf{x}' \rangle},
$$

where $\|\cdot\|$ represents the Euclidean norm \citep{gauche}. 

NNs are a popular choice in QSAR research, though they are often outperformed by more traditional methods \citep{Baskin2008, Koutsoukas2017}. Decision trees and related ensemble models tend to perform stronger on smaller datasets, overfit less, and don't rely on a static 3D molecular structure like NNs do \citep{Baskin2008, Koutsoukas2017}. NNs are often deterministic; however, they can be transformed into probabilistic models through several approaches. BNNs introduce priors on the weights and compute a posterior distribution on those weights given the input data, providing uncertainty quantification. This can be done by approximations, including Monte Carlo dropout, variational inference, or ensembles \citep{gal2016}. To improve computational efficiency, several strategies exist for converting NNs to BNNs. One straightforward yet computationally expensive approach is to replace all linear layers with Bayesian layers (full Bayesian transformation, or full-BNN), using Gaussian priors $\mathcal{N}(0, 0.1^2)$ on all weights. Replacing only the final layer is computationally lighter (last-layer transformation, or last-layer-BNN), using the same priors on the final layer only. A third approach is Variational Bayesian Last Layers (VBLL), which maintains a mean-field variational posterior $q(\mathbf{W}) = \mathcal{N}(\boldsymbol{\mu}_W, \text{diag}(\boldsymbol{\sigma}_W^2))$ over the last-layer weights and is trained by maximizing the evidence lower bound (ELBO), i.e., minimizing the reconstruction loss plus KL divergence $D_{\text{KL}}(q(\mathbf{W}) \| p(\mathbf{W}))$ from a standard normal prior, scaled by $1/N$ \citep{Harrison2024} (VBLL transformation). The VBLL also learns a scalar observation noise variance, enabling decomposition of predictive uncertainty into epistemic (from weight posterior sampling) and aleatoric (from learned noise) components. All BNN variants use 100 Monte Carlo forward passes at inference to estimate predictive distributions. Bayesian approaches also act as forms of regularization \citep{Burden2009}, as dropouts help to prevent overfitting.

\begin{figure}[htbp]
    \centering
    \includegraphics[width=0.85\textwidth]{bayesian_transformation.png}
    \caption{Deterministic versus Bayesian neural network architectures. Deterministic networks use fixed weights $W$ to produce point predictions, while Bayesian networks treat weights as distributions $\mathcal{N}(\mu, \sigma^2)$, generating prediction histograms through multiple stochastic forward passes.}
    \label{fig:bayesian_transformation}
\end{figure}

We used the following deterministic NNs: standard feed-forward NNs and multi-layer perceptron (MLP) networks. To transform these into probabilistic NN variants, we implemented BNNs using three approaches: full Bayesian transformation (full-BNN) with weight uncertainty propagated through all layers, last-layer transformation (last-layer-BNN) where only the final layer is treated probabilistically, and VBLL transformation \citep{Harrison2024}. Full-BNN implementation is visualized in Figure~\ref{fig:bayesian_transformation}. We implemented NN and MLP models using \texttt{PyTorch} \citep{pytorchGeometric} with ReLU activations and dropout regularization ($p=0.2$). The DNN architecture uses two hidden layers of sizes [128, 64] with dropout after each hidden layer. The MLP architecture uses two hidden layers of size [128, 128] with dropout applied before the output layer. Each model was trained using the Adam optimizer with early stopping on validation loss.

\subsection{Performance Metrics}

This paper evaluates performance and noise robustness across a range of metrics. Root mean squared error (RMSE) and the coefficient of determination ($R^2$) are primarily used to compare QSAR performance. $R^2 = 1$ indicates perfect prediction, $R^2 = 0$ indicates performance no better than the mean, and $R^2 < 0$ indicates predictions worse than the mean---which can occur under heavy noise or poor model fits. $\sigma$ denotes the relative noise level for each experiment, standardized across noise strategies. 

Pairwise statistical comparisons between deterministic and probabilistic architecture counterparts (e.g., DNN vs DNN-BNN variants, RF vs QRF) were conducted using the Wilcoxon signed-rank test, with $\alpha = 0.05$ (two-sided). To assess whether model rankings remain consistent across different noise injection strategies, we computed Kendall's coefficient of concordance (Kendall's W), which ranges from 0 (no agreement) to 1 (complete agreement). Values above 0.7 indicate strong agreement in rankings across strategies.

When evaluating uncertainty, we assess the relationship between predicted uncertainty and prediction quality. Spearman's rank correlation coefficient ($\rho$) was used to quantify the relationship between predicted uncertainty ($u_i$) and absolute error ($|y_i - \hat{y}_i|$), as well as between predicted uncertainty and injected noise magnitude. Higher correlations indicate that the model's uncertainty estimates reliably identify predictions that are likely to be incorrect.

For probabilistic models, empirical coverage was evaluated at $1\sigma$ and $2\sigma$ confidence intervals by computing the proportion of test predictions where the true value falls within $\hat{y}_i \pm k\hat{u}_i$:
$$
\text{Coverage}(k\sigma) = \frac{1}{N}\sum_{i=1}^{N} 1\big[|y_i - \hat{y}_i| \leq k\hat{u}_i\big],
$$
such that $\hat{u}_i$ is the predicted uncertainty for sample $i$ and $k \in {1, 2}$. Theoretical target coverages are defined as 68\% at $1\sigma$ and 95\% at $2\sigma$ for Gaussian distributions.

The expected calibration error (ECE) was computed by binning predictions into deciles by predicted uncertainty and measuring the weighted average absolute difference between predicted uncertainty and observed error:
$$
\text{ECE} = \sum_{b=1}^{B} \frac{|B_b|}{N} \big| \bar{u}_b - \bar{e}_b \big|,
$$
where $B_b$ is the set of predictions in bin $b$, $\bar{u}_b$ is the mean predicted uncertainty in that bin, and $\bar{e}_b$ is the mean absolute error. Lower ECE indicates better-calibrated uncertainty estimates.


To evaluate the effect of noise, we examined model performance degradation under increasing noise with $\sigma \in \{0, 0.1, 0.2, \ldots, 1.0\}$ using the Noise Degradation Slope (NDS), computed as the slope of R$^2$ versus noise level:
$$
\text{NDS} = \frac{dR^2}{d\sigma}.
$$

Values closer to zero indicate that noise has minimal effect on performance, negative values indicate a higher sensitivity to noise. Positive values, which would not be observed in this study, would indicate that the noise improves model performance. Configurations with baseline R$^2 \leq 0.6$ were excluded from robustness analysis, as poor performance on clean labels tends to remain poor for noisy labels as well, producing misleadingly shallow slopes.

A two-way analysis of variance (ANOVA) variance decomposition was conducted separately for each noise strategy to identify the relative contributions of molecular representation and model architecture choice on both prediction performance and noise robustness. We chose ANOVA with $\eta^2$ effect sizes rather than pairwise significance tests because our primary question is variance attribution: how much of the variation in robustness is explained by model architecture versus molecular representation. Pairwise tests can establish whether two models differ significantly but cannot partition the total variance among factors---the key distinction our analysis requires. This per-strategy approach avoids inappropriate aggregation across fundamentally different noise types. For a given metric $y$ (either $R^2$ at fixed noise or NDS):
$$
y_{ijr} = \mu + \alpha_i + \beta_j + (\alpha\beta){ij} + \epsilon{ijr},
$$
such that $\alpha_i$ represents the effect of model architecture $i$, $\beta_j$ represents the effect of molecular representation $j$, $(\alpha\beta){ij}$ is the interaction term, and $\epsilon{ijr}$ is the residual for replicate $r$.

The proportion of variance explained by each factor was calculated as the effect size $\eta^2$:
$$
\eta^2_{\text{factor}} = \frac{SS_{\text{factor}}}{SS_{\text{total}}},
$$
using Type I (sequential) sum of squares, such that $SS$ denotes the sum of squares. Both factors (model architecture and molecular representation) are treated as fixed effects, and the 10 experimental replicates per cell provide the error term.

To ensure a valid fully crossed design, we curated the set of ANOVA factor levels by removing redundant or near-duplicate levels. Pairwise Spearman rank correlations were computed between all model NDS profiles and between all representation NDS profiles across noise strategies. Models whose NDS profiles correlated at $\rho > 0.99$ with another model were excluded, including quantile regression forests (QRF), which were highly redundant with RF. Among representations, SNS and the topological fingerprint correlated at $\rho > 0.90$, and SNS was excluded to avoid inflating representation degrees of freedom. Because SNS derives its substructure vocabulary from Morgan circular substructures, and the topological fingerprint encodes related path-based substructures, these two representations capture overlapping structural information. We additionally include a standard Morgan fingerprint (radius 2, 2048 bits) to directly assess whether circular and path-based fingerprints contribute independent signal. We also computed intraclass correlation coefficients ICC(1,1) for all model pairs to assess within-family consistency; ICC(1,1) measures the proportion of total variance attributable to between-subject (i.e., between-configuration) differences, with values near 1.0 indicating that two models rank configurations almost identically. The full redundancy and ICC tables are reported in Supplementary Tables S5--S7. After exclusion, the ANOVA retained all remaining model architectures and six representations (topological fingerprint, Morgan fingerprint, PDV, SMILES, mol2vec, and MHG-GNN), yielding a fully crossed design.

To ensure fair model comparisons, main analyses use a single primary representation (physicochemical descriptor vectors, PDV) rather than aggregating across representations with different baseline performance levels.

\subsection{Noise Strategies}
In this research, our objective is to evaluate the robustness of QSAR models against label noise. To do so, we inject artificial noise into the labels during training. Validation and test data remain free of noise. This artificial noise simulates real-world experimental noise. The clean target label $y_i$ is replaced by a noisy label $\tilde{y}_i = y_i + \epsilon_i$, where $\epsilon_i$ is the artificial noise injected at the index $i$, the value of which is determined by a specified noise strategy.  

Experimental noise is often modeled as homoscedastic Gaussian noise and is added evenly across all labels. However, not all experimental noise is random. It may be heteroscedastic, where the variance depends on the molecule itself or experimental factors, or systematic, where biases from factors like assay conditions are introduced. We used a variety of noise strategies to model the different types of possible experimental noise that could arise in cheminformatics data. 

The following strategies were implemented for regression (Table~\ref{tab:regression_noise}, Figure~\ref{fig:noise_strategies}):
\begin{table}[h]
\centering
\begin{tabular}{lll}
\toprule
\textbf{Noise Type} & \textbf{Noise Scaling} & \textbf{Simulated Real-World Source} \\
\midrule
Legacy & $\sigma$ (fixed) & Random measurement error \\
Outlier & $3\sigma$ / $0.1\sigma$ & Transcription errors, batch effects \\
Quantile & $2\sigma$ / $0.1\sigma$ & Systematic error in hard-to-predict domains \\
Hetero & $\sigma\sqrt{0.1 + 0.05|y|}$ & Heteroscedastic measurement precision \\
Threshold & $2\sigma$ / $0.1\sigma$ & Regime-dependent assay errors \\
Valprop & $\sigma(1 + 0.1|y|)$ & Percentage-based measurement uncertainty \\
\bottomrule
\end{tabular}
\caption{Regression noise injection strategies and their real-world analogues.}
\label{tab:regression_noise}
\end{table}

The parameter $\sigma$ serves as a noise scaling factor across all strategies. For legacy noise, $\epsilon_i \sim \mathcal{N}(0, \sigma^2)$ is applied uniformly. For outlier noise, samples identified as outliers ($z$-score $> 2.0$) receive noise from $\mathcal{N}(0, (3\sigma)^2)$ while normal samples receive $\mathcal{N}(0, (0.1\sigma)^2)$. For quantile noise, samples above the 90th or below the 10th percentile of the target distribution receive noise scaled by $2.0\times$, while central samples receive $0.1\times$. For threshold noise, samples with $|y| > 1.0$ (on normalized data) receive the higher $2.0\times$ multiplier, while those within receive $0.1\times$. For value-proportional, $\epsilon_i \sim \mathcal{N}(0, (\sigma(1 + 0.1|y_i|))^2)$, using a multiplicative formulation where noise scales proportionally with the absolute target value. And finally, for heteroscedastic, noise variance is computed as $\alpha \sigma^2 + \beta \sigma^2 |y_i|$ with $\alpha = 0.1$ and $\beta = 0.05$. All other multipliers remained constant, meaning that while the relative noise distribution differs between strategies, the difficulty scaling controlled by $\sigma$ is consistent. The detailed progression of label corruption across noise levels is shown in Supplementary Figure~\ref{fig:supp_noise_detailed}.

\begin{figure}[htbp]
\centering
\includegraphics[width=\textwidth]{fig_methods_noise_strategies.png}
\caption{Effect of each noise injection strategy on the HOMO-LUMO gap label distribution at $\sigma = 0.5$. Blue outline: clean distribution; colored fill: distribution after noise injection. RMSE between clean and noisy labels is annotated per panel.}
\label{fig:noise_strategies}
\end{figure}

\subsubsection{NoiseInject Framework}

In an extension of this work, we developed an open-source benchmarking framework called NoiseInject to evaluate model robustness and uncertainty quantification in the presence of artificial noise for both regression and classification tasks. The framework accepts datasets in standard formats (CSV, DataFrame, or NumPy arrays) and applies artificial noise to labels while preserving the integrity of the test set.

The framework computes standard regression metrics (RMSE, $R^2$, mean absolute error [MAE]) and classification metrics (accuracy, F1 score, ROC-AUC) along with uncertainty-specific measures: expected calibration error, coverage of the prediction interval (or prediction set for classification) at specified confidence levels, mean interval width, and Pearson correlation between predicted uncertainty and absolute error. We included noise robustness metrics including noise degradation slope and retention percentage to quantify performance degradation rates. All experiments are logged to JSON Lines format with full metadata (noise configuration, model identifiers, hyperparameters, random seeds).

The evaluation pipeline is model-agnostic, although reference implementations of GPs (using \texttt{Gauche} \citep{gauche}), split-CP for sklearn-compatible models, and Monte Carlo dropout wrappers for PyTorch networks are provided. We added visualization tools which generate performance degradation curves, calibration plots, and uncertainty-error scatter plots for rapid assessment of model behavior under noise. 

The software is available under an MIT license at \url{github.com/adpunt/noise\_inject} and is designed to integrate into existing cheminformatics or other data pipelines, enabling systematic benchmarking of model robustness.


\section{Results}

\subsection{The Role Inversion: Performance vs Robustness}\label{sec:role_inversion}

A central finding of this study is that the factors governing predictive performance differ fundamentally from those governing noise robustness. ANOVA decomposition across six noise strategies reveals a consistent pattern (Table~\ref{tab:anova_decomposition}, Figure~\ref{fig:anova_decomposition}). For predictive performance (R$^2$ at $\sigma = 0.3$), the model--representation interaction term is the single largest source of variance across all six strategies, followed by representation and then model architecture. This means that specific model--representation \emph{pairings} matter more than either factor alone: a model's performance depends on which representation it receives, and vice versa.

For noise robustness (NDS), this ordering inverts. Model architecture becomes the dominant factor, the interaction term drops to second, and representation recedes to a minor role. We call this the \emph{role inversion}: the factors that govern how well a model predicts are not the same factors that govern how well it resists noise.

The role inversion has a natural interpretation. For performance, molecular representation determines what chemical information is available to the model, while architecture determines how efficiently it is used---the large interaction term reflects the fact that these two choices are not independent. Certain models extract more from certain representations, and the best-performing configurations are specific pairings rather than universally superior models or representations. For robustness, representation becomes less important because noise corrupts labels, not features. The model's inductive biases---how it regularizes, ensembles, or distributes uncertainty---determine whether corrupted labels degrade predictions or are absorbed by the learning algorithm. Model-specific properties dominate regardless of which representation is used.

The strength of these effects depends on the noise strategy (Figure~\ref{fig:anova_decomposition}B). The strongest model effects emerge under the most severe strategies---threshold and value-proportional noise---where deterministic corruption of specific label regions forces architectures to reveal their inductive biases. Under outlier noise, which affects only statistical outliers, model and interaction effects are both small ($\eta^2_{\text{model}} = 12\%$, $\eta^2_{\text{interaction}} = 6\%$) while the residual dominates ($80\%$): when noise is mild and targeted, individual variation overwhelms systematic factors. At the other extreme, under value-proportional noise, where noise scales with the target value itself, model architecture alone explains over half the variance ($\eta^2_{\text{model}} = 55\%$)---the noise is severe enough that only architectural properties can protect against it. To our knowledge, this variance decomposition of noise robustness across both model and representation factors has not been reported previously in the cheminformatics literature.

\begin{table}[htbp]
\centering
\caption{ANOVA variance decomposition by noise strategy. $\eta^2$ (\%) for model architecture and molecular representation effects on performance (R$^2$ at $\sigma=0.3$) and robustness (NDS).}
\label{tab:anova_decomposition}
\small
\begin{tabular}{lcccccccc}
\toprule
& \multicolumn{4}{c}{\textbf{Performance ($\eta^2$, \%)}} & \multicolumn{4}{c}{\textbf{Robustness ($\eta^2$, \%)}} \\
\cmidrule(lr){2-5} \cmidrule(lr){6-9}
\textbf{Strategy} & \textbf{Model} & \textbf{Rep} & \textbf{Inter.} & \textbf{Resid.} & \textbf{Model} & \textbf{Rep} & \textbf{Inter.} & \textbf{Resid.} \\
\midrule
Gaussian        & 26.8 & 25.2 & 44.2 &  3.9 & 43.9 &  7.0 & 21.9 & 27.2 \\
Quantile        & 27.8 & 25.4 & 42.7 &  4.0 & 45.0 &  4.6 & 19.1 & 31.3 \\
Threshold       & 29.3 & 24.7 & 42.1 &  3.9 & 48.8 &  8.2 & 34.2 &  8.8 \\
Heteroscedastic & 28.0 & 26.7 & 42.7 &  2.6 & 39.6 &  3.2 & 20.6 & 36.6 \\
Value-prop.     & 28.7 & 24.7 & 42.8 &  3.7 & 54.8 &  6.5 & 28.0 & 10.8 \\
Outlier         & 26.5 & 28.6 & 44.3 &  0.6 & 12.2 &  1.0 &  6.3 & 80.4 \\
\bottomrule
\end{tabular}
\end{table}

\begin{figure}[htbp]
    \centering
    \includegraphics[width=\textwidth]{fig2_anova_decomposition.png}
    \caption{ANOVA variance decomposition ($\eta^2$, \%) for (A) predictive performance (R$^2$ at $\sigma=0.3$) and (B) noise robustness (NDS), by noise strategy.}
    \label{fig:anova_decomposition}
\end{figure}

The interaction between model architecture and representation is visualized directly in Figure~\ref{fig:interaction}. Panel~A reveals the structure of this interaction: some models (SVM, full BNNs) maintain consistent robustness across all representations, producing uniform rows in the heatmap, while others (MLP, last-layer BNNs) show steep variation across representations, indicating that their robustness depends critically on the input representation. Panel~B confirms that model robustness on one representation partially predicts robustness on another (Spearman $\rho = 0.67$, $p = 0.024$ for topological fingerprint vs PDV), but imperfectly---the scatter around the diagonal reflects the interaction term identified by the ANOVA. These patterns foreshadow the three distinct robustness mechanisms analyzed in Section~\ref{sec:mechanisms}.

\begin{figure}[htbp]
\centering
\includegraphics[width=\textwidth]{fig_interaction.png}
\caption{Representation--model interaction effects. (A) NDS heatmap showing model robustness across representations under Gaussian noise. (B) NDS on PDV versus topological fingerprint for each model under Gaussian noise, with Spearman correlation.}
\label{fig:interaction}
\end{figure}

To ensure meaningful robustness comparisons, we restricted analysis to configurations achieving baseline R$^2 > 0.6$, excluding 95 configurations where poor clean-data performance would produce misleadingly shallow degradation slopes (Supplementary Table~\ref{tab:excluded_configs}).

\subsection{Model Rankings and Noise Strategy Patterns}\label{sec:rankings}

We evaluated all model architectures under increasing noise on the PDV representation across six strategies. Performance degrades approximately linearly with noise level for all models, but the rate of degradation---quantified by the Noise Degradation Slope (NDS)---varies substantially across architectures (Figure~\ref{fig:global_overview}A). The NDS heatmap (Figure~\ref{fig:global_overview}B) reveals that these differences persist across all six noise strategies, with threshold and value-proportional noise producing the steepest degradation for all models.

\begin{figure}[htbp]
\centering
\includegraphics[width=\textwidth]{fig1_global_overview.png}
\caption{Global overview of noise robustness. (A) Performance degradation curves (PDV, Gaussian noise) showing R$^2$ versus noise level $\sigma$ for representative models. (B) NDS heatmap across all model architectures and noise strategies on PDV.}
\label{fig:global_overview}
\end{figure}

NGBoost and SVM showed the smallest degradation across all strategies, while MLP and last-layer BNNs showed the steepest (Table~\ref{tab:nds_ranking}). Model rankings across noise strategies were highly consistent: Kendall's $W = 0.949$ ($p = 1.5 \times 10^{-10}$; 14 models, 6 strategies), indicating that a model's relative robustness is largely strategy-independent (Supplementary Figure~\ref{fig:supp_ranking_consistency}). Parallel analysis on the topological fingerprint confirmed the same patterns (Supplementary Figures~\ref{fig:supp_ecfp4_overview} and~\ref{fig:supp_ecfp4_ranking}).

\begin{table}[htbp]
\centering
\caption{NDS by model on PDV, ranked by mean across six strategies. Lower $|$NDS$|$ indicates greater robustness. Full results in Supplementary Table~\ref{tab:nds_all_reps}.}
\label{tab:nds_ranking}
\small
\begin{tabular}{lccccccc}
\toprule
\textbf{Model} & \textbf{Gauss.} & \textbf{Outlier} & \textbf{Quant.} & \textbf{Thresh.} & \textbf{Hetero.} & \textbf{Val.-P.} & \textbf{Mean} \\
\midrule
NGBoost             & $-0.31$ & $-0.10$ & $-0.27$ & $-0.61$ & $-0.17$ & $-0.54$ & $-0.33$ \\
SVM                 & $-0.34$ & $-0.10$ & $-0.28$ & $-0.67$ & $-0.18$ & $-0.57$ & $-0.36$ \\
MLP-BNN (Full)      & $-0.35$ & $-0.12$ & $-0.30$ & $-0.66$ & $-0.19$ & $-0.59$ & $-0.37$ \\
DNN-BNN (Full)      & $-0.34$ & $-0.15$ & $-0.30$ & $-0.66$ & $-0.19$ & $-0.59$ & $-0.37$ \\
MLP-VBLL      & $-0.35$ & $-0.12$ & $-0.30$ & $-0.67$ & $-0.20$ & $-0.60$ & $-0.37$ \\
LightGBM            & $-0.35$ & $-0.11$ & $-0.30$ & $-0.69$ & $-0.19$ & $-0.61$ & $-0.38$ \\
DNN-VBLL      & $-0.36$ & $-0.12$ & $-0.30$ & $-0.69$ & $-0.20$ & $-0.61$ & $-0.38$ \\
XGBoost             & $-0.35$ & $-0.12$ & $-0.31$ & $-0.69$ & $-0.19$ & $-0.62$ & $-0.38$ \\
GP (Gauche)         & $-0.37$ & $-0.12$ & $-0.32$ & $-0.69$ & $-0.20$ & $-0.62$ & $-0.39$ \\
DNN                 & $-0.36$ & $-0.13$ & $-0.33$ & $-0.70$ & $-0.20$ & $-0.63$ & $-0.39$ \\
RF                  & $-0.36$ & $-0.12$ & $-0.31$ & $-0.73$ & $-0.20$ & $-0.64$ & $-0.39$ \\
DNN-BNN (Last)      & $-0.36$ & $-0.13$ & $-0.33$ & $-0.69$ & $-0.23$ & $-0.62$ & $-0.39$ \\
MLP                 & $-0.36$ & $-0.13$ & $-0.34$ & $-0.72$ & $-0.21$ & $-0.65$ & $-0.40$ \\
MLP-BNN (Last)      & $-0.39$ & $-0.14$ & $-0.34$ & $-0.70$ & $-0.22$ & $-0.64$ & $-0.41$ \\
\bottomrule
\end{tabular}
\end{table}

The six strategies fell into three severity tiers (Table~\ref{tab:phase5_noise_strategy}): outlier and heteroscedastic noise were mild ($>93\%$ performance retained), quantile and Gaussian were moderate, and threshold and value-proportional were severe ($>34\%$ performance lost). Threshold noise was the most destructive because it deterministically corrupts all labels above a property cutoff, introducing systematic bias that models cannot learn around.

\begin{table}[htbp]
\centering
\caption{Noise strategy severity: mean performance retention and NDS across ANOVA-included configurations.}
\label{tab:phase5_noise_strategy}
\small
\begin{tabular}{lrrrr}
\toprule
\textbf{Noise} & \textbf{Mean} & \textbf{Mean} & \textbf{Std} & \textbf{Mean} \\
\textbf{Strategy} & \textbf{Baseline $R^2$} & \textbf{Retention \%} & \textbf{Retention \%} & \textbf{$|$NDS$|$} \\
\midrule
Outlier & 0.778 & 97.5 & 0.3 & 0.099 \\
Heteroscedastic & 0.776 & 93.4 & 0.7 & 0.194 \\
Quantile & 0.785 & 88.6 & 1.5 & 0.286 \\
Gaussian & 0.831 & 81.0 & 3.3 & 0.362 \\
Value-proportional & 0.785 & 65.9 & 3.6 & 0.612 \\
Threshold & 0.785 & 57.8 & 4.7 & 0.675 \\
\bottomrule
\end{tabular}
\end{table}

Despite the variation in absolute NDS values, the \emph{relative} severity of noise strategies was remarkably uniform across architectures: sensitivity ratios (each model's NDS on a given strategy divided by its Gaussian NDS) varied by less than $10\%$ across all models for every strategy (Supplementary Table~\ref{tab:supp_sensitivity_ratios}), indicating that strategy severity is a property of the noise structure, not the model.

Baseline performance does not predict noise robustness (Supplementary Figure~\ref{fig:supp_full_overview}). Models with inherent noise robustness tend to sacrifice clean-data accuracy---NGBoost and SVM rank lower on clean data but rise under heavy noise (Supplementary Table~\ref{tab:sigma_rankings})---a trade-off that only becomes apparent under systematic noise evaluation.

DNN variants with different hidden layer configurations ([128,64], [256,128,64], [512,256]) produced nearly identical NDS values, confirming that noise resistance is driven by architectural class rather than network capacity.

\subsection{Three Mechanisms of Noise Robustness}\label{sec:mechanisms}

The overall ANOVA establishes that model architecture is the dominant factor for robustness, but simple effects analysis (Supplementary Table~\ref{tab:simple_effects}) reveals that different models achieve robustness through qualitatively different mechanisms, depending on how much their noise robustness depends on the choice of representation.

\paragraph{Inherent robustness.} SVM and full BNNs show robustness that is largely representation-independent: within these models, representation explains less than 15\% of robustness variance across all strategies (Supplementary Table~\ref{tab:simple_effects}). These models achieve robustness through their inductive biases---SVM's margin maximization and the BNN's weight priors---rather than by relying on a favorable input representation.

\paragraph{Representation-mediated robustness.} RF, MLP, and last-layer BNNs show the opposite pattern. Within MLP, representation explains over 60\% of robustness variance under Gaussian noise and over 85\% under threshold noise. RF shows similarly strong representation dependence (Supplementary Table~\ref{tab:simple_effects}). These models can be robust, but only when paired with the right representation---on the wrong representation, they degrade steeply.

\paragraph{The BNN natural experiment.} Full versus last-layer Bayesian transformation provides a controlled comparison: same base architecture, same training data, same representations, differing only in the extent of Bayesian treatment. The contrast is stark. For DNN-BNN-Last, representation explains $57.7\%$ of robustness variance under Gaussian noise; for DNN-BNN-Full, only $15.3\%$. For MLP variants, the difference is even larger: $79.6\%$ (Last) versus $28.0\%$ (Full). Full Bayesian treatment transforms a representation-mediated architecture into an inherently robust one by placing priors on all weights, not just the output layer.

However, VBLL's robustness advantage is representation-dependent. On MHG-GNN and mol2vec representations, VBLL models were catastrophically unstable, with R$^2$ frequently below $-100$ and only 1--4 usable iterations out of 10. All 24 VBLL $\times$ \{MHG-GNN, mol2vec\} configurations were excluded due to baseline R$^2 < 0.6$. By contrast, full BNNs trained successfully on all representations, reinforcing the distinction between inherent robustness (full BNN) and representation-mediated robustness (VBLL).

This mechanistic difference is visible in the degradation curves (Figures~\ref{fig:dnn_family} and~\ref{fig:mlp_rf_comparison}). For the DNN family on PDV, full BNN transformation reduces the Gaussian NDS from $-0.363$ (DNN) to $-0.340$ (DNN-BNN Full), while last-layer transformation provides no improvement ($-0.358$; Table~\ref{tab:wilcoxon_bnn}). VBLL falls between the two ($-0.359$), achieving comparable robustness to the full BNN with less computational cost. The pattern is more pronounced under heteroscedastic noise, where DNN-BNN (Last) degrades more steeply ($-0.226$) than both the full BNN ($-0.188$) and VBLL ($-0.197$). Full BNN transformation significantly improved robustness for both DNN ($p < 10^{-6}$) and MLP ($p < 10^{-7}$), while last-layer transformation provided no significant benefit. VBLL transformation also significantly improved robustness for both architectures ($p < 10^{-4}$), with MLP-VBLL achieving the largest improvement ($\Delta \text{NDS} = +0.134$; Table~\ref{tab:wilcoxon_bnn}). The benefit of full transformation varies by representation (Supplementary Table~\ref{tab:supp_nn_transforms}), with the largest improvement on SMILES where the base model struggles most.

\begin{table}[htbp]
\centering
\caption{Wilcoxon signed-rank tests for Bayesian and probabilistic transformations. Positive $\Delta$ NDS indicates the variant is more robust than its deterministic counterpart.}
\label{tab:wilcoxon_bnn}
\small
\begin{tabular}{llrrrl}
\toprule
\textbf{Family} & \textbf{Comparison} & \textbf{$n$} & \textbf{$\Delta$ NDS} & \textbf{$p$-value} & \textbf{Sig.} \\
\midrule
DNN & DNN vs DNN-BNN (Full) & 24 & $+0.066$ & $6.0 \times 10^{-7}$ & * \\
DNN & DNN vs DNN-BNN (Last) & 22 & $+0.002$ & $0.17$ & \\
DNN & DNN vs DNN-VBLL & 18 & $+0.069$ & $5.3 \times 10^{-5}$ & * \\
MLP & MLP vs MLP-BNN (Full) & 24 & $+0.115$ & $1.2 \times 10^{-7}$ & * \\
MLP & MLP vs MLP-BNN (Last) & 22 & $+0.002$ & $0.42$ & \\
MLP & MLP vs MLP-VBLL & 18 & $+0.134$ & $7.6 \times 10^{-6}$ & * \\
RF  & RF vs QRF             & 34 & $-0.023$ & $1.6 \times 10^{-9}$ & * \\
\bottomrule
\end{tabular}
\end{table}

For random forests, QRF was significantly \emph{less} robust than RF (Wilcoxon $p = 1.6 \times 10^{-9}$, mean NDS worsening $-0.023$). The QRF's quantile regression mechanism does not confer noise robustness; instead, the additional flexibility may cause overfitting to noisy labels.

\begin{figure}[htbp]
    \centering
    \includegraphics[width=\textwidth]{fig4_dnn_family.png}
    \caption{DNN family: R$^2$ versus $\sigma$ for deterministic DNN and Bayesian variants (full BNN, last-layer BNN, VBLL) under Gaussian and heteroscedastic noise (PDV).}
    \label{fig:dnn_family}
\end{figure}

\begin{figure}[htbp]
    \centering
    \includegraphics[width=\textwidth]{fig5_mlp_rf_comparison.png}
    \caption{MLP and RF families: R$^2$ versus $\sigma$ for (A--B) MLP and Bayesian variants (full BNN, last-layer BNN, VBLL) and (C--D) RF versus QRF, under Gaussian and heteroscedastic noise (PDV).}
    \label{fig:mlp_rf_comparison}
\end{figure}

The representation perspective complements the model perspective. Some representations amplify architectural differences---on SMILES, model choice explains over 90\% of robustness variance---while others compress them: on PDV, model explains only 40\% (Supplementary Table~\ref{tab:simple_effects}). PDV acts as an equalizer, producing moderate robustness regardless of model choice, while SMILES acts as an amplifier, making model selection critical (Figure~\ref{fig:interaction}; Supplementary Figure~\ref{fig:supp_full_overview}; Supplementary Table~\ref{tab:nds_all_reps}).

\subsection{Uncertainty as the Mechanistic Link}\label{sec:uncertainty}

The three robustness mechanisms suggest a deeper question: \emph{why} do some architectures resist noise? We hypothesized that models whose uncertainty estimates respond to injected noise would also maintain prediction quality under noise---that noise detection and noise resistance are linked. Among the nine uncertainty-producing models (Table~\ref{tab:uncertainty_metrics}), models with explicit noise-modeling channels (NGBoost, GP) show the strongest noise detection (Unc-Noise $\rho = 0.36$--$0.40$), followed by full BNNs and QRF at moderate levels ($\rho = 0.22$--$0.25$), with last-layer BNNs and VBLL variants at the bottom ($\rho = 0.10$--$0.14$). The VBLL variants are a notable exception to the noise-detection/robustness link---achieving robustness comparable to full BNNs despite weak noise tracking---suggesting that their learned observation noise parameter absorbs noise during training even when the resulting uncertainty estimates are less responsive at test time.

\begin{table}[htbp]
\centering
\caption{Uncertainty quantification metrics for probabilistic models (legacy noise strategy, averaged across representations). Unc-Error $\rho$: Spearman correlation between predicted uncertainty and absolute prediction error. Unc-Noise $\rho$: Spearman correlation between predicted uncertainty and injected noise level $\sigma$. Coverage targets: 68\% ($1\sigma$), 95\% ($2\sigma$). Per-representation and per-strategy breakdowns are in Supplementary Table~\ref{tab:supp_uncertainty_by_strategy_rep}.}
\label{tab:uncertainty_metrics}
\small
\begin{tabular}{lrrrrr}
\toprule
\textbf{Model} & \textbf{Unc-Error $\rho$} & \textbf{Unc-Noise $\rho$} & \textbf{ECE} & \textbf{Cov. $1\sigma$} & \textbf{Cov. $2\sigma$} \\
\midrule
QRF            & 0.22 & 0.25 & 0.16 & 69.9\% & 92.4\% \\
NGBoost        & 0.22 & 0.40 & 0.12 & 69.4\% & 95.0\% \\
GP (Gauche)    & 0.20 & 0.36 & 0.19 & 74.7\% & 95.0\% \\
MLP-BNN (Last) & 0.13 & 0.14 & 0.14 & 51.5\% & 78.7\% \\
MLP-BNN (Full) & 0.13 & 0.22 & 0.41 & 66.1\% & 88.4\% \\
DNN-BNN (Last) & 0.09 & 0.13 & 0.18 & 47.5\% & 74.0\% \\
DNN-BNN (Full) & 0.04 & 0.24 & 0.20 & 62.3\% & 85.4\% \\
MLP-VBLL & $-0.04$ & 0.10 & 8.02 & 59.3\% & 77.8\% \\
DNN-VBLL & $-0.08$ & 0.10 & 4.98 & 56.3\% & 74.4\% \\
\bottomrule
\end{tabular}
\end{table}

\begin{figure}[htbp]
\centering
\includegraphics[width=\textwidth]{fig_uncertainty_combined.png}
\caption{(A) Mean predicted uncertainty versus $\sigma$ for each probabilistic model (PDV, Gaussian strategy). (B) Aleatoric and epistemic uncertainty components versus $\sigma$ for models with decomposition data.}
\label{fig:uncertainty_combined}
\end{figure}

NGBoost exemplifies this link: it produces the strongest noise tracking and achieves both the best robustness and calibration (Table~\ref{tab:uncertainty_metrics}). Its advantage is mechanistic: NGBoost minimizes the negative log-likelihood of a learned Normal distribution, providing two gradient pathways when a noisy label produces a large residual---the model can increase the predicted scale $\sigma(x)$ rather than distorting the mean $\mu(x)$. The loss function routes noise into the uncertainty estimate instead of into the prediction. At the other extreme, QRF derives uncertainty from post-hoc quantile spreads across trees, with no training-time mechanism to respond to label quality---its noise tracking is weak ($\rho = 0.25$) relative to models with explicit noise channels, mapping onto its poor robustness.

The full versus last-layer BNN pattern reinforces this from the uncertainty side (Table~\ref{tab:uncertainty_metrics}): full variants show moderate noise tracking ($\rho = 0.22$--$0.24$) and higher coverage than last-layer variants, though both undercover relative to theoretical targets. The same architectural change that converts representation-mediated robustness into inherent robustness (Section~\ref{sec:mechanisms}) also converts noise-blind uncertainty into noise-responsive uncertainty---the two properties are mechanistically linked through the extent of Bayesian treatment (Figure~\ref{fig:uncertainty_combined}).

The critical architectural distinction is whether a model's loss function provides an explicit gradient pathway for noise absorption. Models trained with likelihood-based objectives---NGBoost, GP, and VBLL---contain a learned variance or noise parameter that creates a trade-off: when a noisy label produces a large residual, increasing the noise parameter reduces the likelihood penalty at lower cost than distorting the mean prediction. Models trained with squared-error objectives---standard BNNs---lack this trade-off. All training pressure goes toward fitting the mean, and the weight posterior broadens only as a side effect of conflicting gradients from noisy labels, not as a directed response.

Uncertainty decomposition makes this distinction visible (Figure~\ref{fig:uncertainty_combined}B). In the GP (Gauche), which separates posterior variance from a learned observation noise term, the observation noise component rises steeply with $\sigma$ while posterior variance remains flat---the model routes injected noise into its dedicated noise parameter rather than absorbing it into the prediction. BNNs, which lack a noise parameter, respond to noise entirely through posterior broadening, producing weaker noise tracking ($\rho = 0.22$--$0.24$ versus $\rho > 0.36$ for models with explicit noise channels). The VBLL variants add a learned observation noise parameter to the neural network loss, providing the same active noise channel within a BNN architecture and bridging this gap. However, the VBLL results reveal an important subtlety: despite achieving robustness comparable to full BNNs (Table~\ref{tab:nds_ranking}), VBLL models show weak noise tracking ($\rho = 0.10$, Table~\ref{tab:uncertainty_metrics}) and catastrophically poor calibration (ECE $= 4.98$--$8.02$), indicating that while the learned noise variance absorbs noise during training, it does not produce well-calibrated uncertainty at inference. This dissociation suggests that noise absorption during training and noise responsiveness at inference are partially separable. VBLL's learned noise variance acts as a global regularizer that prevents overfitting to noisy labels, but because it is shared across all instances rather than conditioned on input, the resulting uncertainty estimates are less instance-specific than those of NGBoost or GP. The practical implication is that VBLL provides robustness without reliable per-prediction noise detection---a useful property when robustness matters more than calibrated uncertainty.

These uncertainty patterns are not specific to a single noise type or representation: the uncertainty--noise correlation structure holds across all six strategies and four representations (Supplementary Table~\ref{tab:supp_uncertainty_by_strategy_rep}). The practical implication is that only models whose loss functions provide an explicit channel for noise absorption---allowing the model to widen its predictive distribution rather than chase noisy labels---provide the kind of noise resistance that generalizes.

\subsection{Generalization Across Strategies and Datasets}\label{sec:generalization}

The uniform strategy sensitivity ratios (Section~\ref{sec:rankings}) establish that robustness findings generalize well across noise types: a model's relative robustness on one strategy predicts its robustness on others ($W = 0.949$). The noise strategy acts as a scaling factor, not a differentiator---it determines how much damage occurs, but not which models handle it best.

Generalization across datasets is a different matter. We evaluated seven representative models---NGBoost, SVM, RF, QRF, DNN, LightGBM, and XGBoost---on three experimentally-derived molecular property datasets: LogD (lipophilicity), Caco-2 efflux permeability, and hERG-Ki (cardiac ion channel binding affinity), using all six noise strategies at $\sigma \in \{0, 0.1, \ldots, 1.0\}$. BNN variants and GP (Gauche) were not included in external validation due to computational constraints. External validation used 5-fold scaffold cross-validation (GroupKFold on Murcko scaffolds) rather than the single deterministic scaffold split used for QM9, providing fold-averaged robustness estimates.

A lower baseline R$^2$ threshold of 0.3 (versus 0.6 for QM9) was applied, reflecting the inherently weaker baselines on smaller experimental datasets: LogD baselines ranged from $R^2 = 0.63$--$0.82$, hERG-Ki from $0.40$--$0.56$, and Caco-2 from $0.37$--$0.62$. PDV representations failed catastrophically on these datasets---DNN+PDV on hERG-Ki produced $R^2 = -2096$, and several other model+PDV configurations yielded negative $R^2$---suggesting that quantum-chemical descriptors optimized for QM9 do not transfer to drug-like property endpoints. Configurations with $|$NDS$| > 2$ were additionally filtered as artifacts (``N/A'' in Figure~\ref{fig:validation_overview}).

\begin{figure*}[htbp]
\centering
\includegraphics[width=\textwidth]{fig_validation_overview.png}
\caption{NDS heatmaps for three external validation datasets (LogD, Caco-2, hERG-Ki). Black cells indicate missing configurations; ``N/A'' indicates filtered extreme values ($|$NDS$| > 2$).}
\label{fig:validation_overview}
\end{figure*}

ANOVA on the external datasets confirms the role inversion: model architecture dominates robustness variance on all three datasets ($\eta^2 = 91$--$95\%$ for model versus $0.7$--$1.8\%$ for representation; Supplementary Figure~\ref{fig:supp_validation_anova}, Table~\ref{tab:supp_validation_anova}). The effect is strongest on Caco-2 ($\eta^2_{\text{model}} = 95.4\%$) and similarly strong on hERG-Ki ($91.6\%$) and LogD ($90.9\%$). Despite this structural consistency, validation rankings diverge substantially from QM9 (Figure~\ref{fig:validation_model_comparison}). NGBoost achieved the best mean rank across external datasets, followed by RF and SVM, while XGBoost performed worst. The correlation between QM9 and external dataset robustness is weak and non-significant ($r = 0.14$, $p = 0.77$; Figure~\ref{fig:validation_qm9_transferability}), indicating that overall robustness rankings do not transfer.

\begin{figure}[htbp]
\centering
\includegraphics[width=0.8\textwidth]{fig_validation_qm9_transferability.png}
\caption{Robustness transferability from QM9 to external datasets. Each point is a model; axes show mean NDS on QM9 versus mean NDS across three external datasets. The weak, non-significant correlation ($r = 0.14$, $p = 0.77$) indicates that robustness rankings do not transfer across domains. NGBoost (noise-channeling robustness) performs well on both; XGBoost (dataset-specific robustness) collapses externally.}
\label{fig:validation_qm9_transferability}
\end{figure}

This non-transferability is interpretable through the mechanism framework developed in Section~\ref{sec:mechanisms}. NGBoost, whose robustness derives from its noise-channeling loss function, ranks first on both QM9 and external datasets---its mechanism of routing noise into the predictive distribution transfers across data domains. SVM, whose robustness is inherent and representation-independent, also transfers well (rank 3). XGBoost, whose QM9 robustness benefits from the large dataset size (130k molecules) and well-separated scaffold structure, loses this advantage on smaller, experimentally-derived datasets with 1,500--5,000 molecules. The mechanisms that produce robustness, not just the robustness rankings themselves, determine whether findings transfer.

\begin{figure*}[htbp]
\centering
\includegraphics[width=\textwidth]{fig_validation_model_comparison.png}
\caption{Model robustness comparison across QM9 and three external validation datasets. Mean NDS by model for each dataset; models ordered by overall mean NDS (most robust first). QM9 rankings diverge substantially from external datasets, particularly for gradient boosting models.}
\label{fig:validation_model_comparison}
\end{figure*}

The gradient boosting family illustrates this starkly. On QM9, NGBoost ranks first in robustness thanks to its noise-channeling loss function, while XGBoost and LightGBM rank in the middle of the pack. On external datasets, XGBoost collapses on Caco-2 efflux (NDS $\approx -1.2$), while LightGBM degrades half as steeply ($\text{NDS} \approx -0.7$). XGBoost's greedy exact splits overfit to QM9's large, well-structured training set; on smaller experimental datasets, this learned structure becomes a liability. LightGBM's histogram-based binning provides modest regularization but does not prevent the same failure mode. NGBoost's robustness transfers fully to external datasets, where it again ranks first---its probabilistic framework, which learns a full predictive distribution rather than a point estimate, provides built-in protection: when label quality degrades, the model widens its predictive intervals rather than producing overconfident wrong predictions.

Certain findings were nonetheless robust across datasets. QRF was consistently less robust than RF on every external dataset (Supplementary Table~\ref{tab:supp_validation_probabilistic}), confirming the QM9 finding. Caco-2 efflux was the most challenging dataset, with all models showing steeper degradation than on LogD or hERG-Ki (Supplementary Figure~\ref{fig:supp_validation_strategy}). Representation choice had minimal effect on robustness across all three external datasets (Figure~\ref{fig:validation_rep_comparison}), reinforcing the ANOVA finding that model architecture, not representation, drives noise robustness.

\begin{figure}[htbp]
\centering
\includegraphics[width=\textwidth]{fig_validation_rep_comparison.png}
\caption{Representation effect on robustness across validation datasets. Mean NDS by representation and dataset, confirming that representation has minimal influence on noise robustness compared to model architecture.}
\label{fig:validation_rep_comparison}
\end{figure}

\section{Conclusion}\label{sec13}
This study demonstrates that while model architecture and molecular representation contribute roughly equally to QSAR model performance, model architecture becomes the dominant factor under label noise---explaining up to 55\% of robustness variance compared to at most 8\% for representation. To our knowledge, this is the first study to combine systematic noise injection across multiple representations and model architectures, uncertainty strategy comparison, and external validation within a single framework. The variance decomposition approach, unusual in cheminformatics where pairwise significance tests predominate, enables a new kind of question: not just ``which model is better,'' but ``what explains the variation in robustness.''

Simple effects analysis reveals three distinct mechanisms of noise robustness: \emph{inherent} robustness (SVM, full BNNs) that operates independently of representation choice, \emph{representation-mediated} robustness (RF, MLP, last-layer BNNs) that depends critically on the input features, and ensemble-based robustness (LightGBM, XGBoost) that benefits from dataset-specific structure. The BNN natural experiment demonstrates that full Bayesian treatment converts a representation-mediated architecture into an inherently robust one ($p < 10^{-6}$), while last-layer treatment provides no benefit. Variational Bayesian Last Layers (VBLL) similarly improve robustness ($p < 10^{-4}$), offering a computationally lighter path to Bayesian regularization---though with a trade-off: VBLL achieves robustness comparable to full BNNs but produces less responsive uncertainty estimates, providing noise absorption during training without reliable per-prediction noise detection at inference.

These mechanisms are linked through uncertainty quantification: among probabilistic models, those with explicit noise-modeling channels (NGBoost, GP, full BNNs) best detect noise and best resist it. The critical architectural distinction is whether a model's loss function provides an explicit gradient pathway for noise absorption---allowing the model to widen its predictive distribution rather than chase noisy labels. NGBoost exemplifies this: it produces the best-calibrated intervals and strongest noise tracking among all models tested. QRF, whose uncertainty responds only modestly to noise, is significantly less robust than its deterministic RF counterpart across all datasets tested.

Across all six noise strategies, model rankings were highly concordant ($W = 0.949$), and the relative severity of noise strategies was uniform across architectures---the noise structure, not the model, determines relative damage. However, robustness rankings did not transfer from QM9 to smaller experimental datasets ($r = 0.14$, $p = 0.77$). Critically, whether findings transfer depends on the \emph{mechanism}: NGBoost's noise-channeling robustness and SVM's inherent robustness generalize to new domains, while XGBoost's dataset-specific robustness does not.

These findings yield practical guidance for model selection under noisy conditions. When robustness must transfer across domains---as is common when moving from large benchmarks to smaller experimental assays---SVM and full BNNs are the safest choices, as their noise resistance is representation-independent and mechanism-driven. When calibrated uncertainty is required, NGBoost provides the strongest noise tracking and best calibration; VBLL provides comparable robustness at lower computational cost, but its uncertainty estimates are less informative about per-prediction noise. When a favorable representation is available and computational cost is a concern, simpler models such as RF can achieve adequate robustness, provided the representation is well-matched to the task. Importantly, when the noise type is unknown, the concordant rankings mean that a model's robustness on one noise type predicts its robustness on others---there is no need to match the noise strategy to the suspected noise source.

Several limitations should be acknowledged. The primary experiments were conducted on QM9, a large ($N = 130{,}000$), computationally derived dataset with negligible measurement noise. While the external validation on three experimentally-derived ADME datasets provides initial evidence for generalizability, these datasets are small ($N = 1{,}482$--$5{,}039$) and represent only a narrow slice of the chemical and biological space relevant to drug discovery. Only seven of the fourteen ANOVA-included models were evaluated externally due to computational constraints, and probabilistic models such as GP (Gauche) and BNN variants were not included in external validation. VBLL models showed training instability on certain representations (mhggnn, mol2vec), with frequent catastrophic iterations that required filtering; this limits the generality of the VBLL findings to representations where training is stable. Additionally, all experiments used regression targets; classification tasks, which are common in ADME and toxicity prediction, may exhibit different robustness patterns.

Future work should extend this framework to classification tasks, larger and more diverse experimental datasets, and additional probabilistic architectures including deep ensembles and evidential neural networks. The VBLL instability on certain representations warrants investigation into training stabilization techniques. Domain-specific noise models---for example, incorporating known assay variability patterns---could replace the generic noise strategies used here and provide more realistic benchmarks for specific therapeutic areas.

To facilitate adoption of noise robustness evaluation as a standard component of QSAR model development, we release NoiseInject, an open-source, model-agnostic benchmarking framework. NoiseInject supports both regression and classification tasks, implements all six noise strategies described here, and computes the full suite of robustness and uncertainty metrics (NDS, coverage, ECE, uncertainty--error correlation) through a unified API. The framework accepts arbitrary datasets and models, enabling researchers to benchmark new architectures and representations against the baselines established in this study. NoiseInject is available at \url{github.com/adpunt/noise\_inject} under an MIT license.

\section*{Declarations}

\paragraph{Competing interests}
The authors declare no competing interests.

% Some journals require declarations to be submitted in a standardised format. Please check the Instructions for Authors of the journal to which you are submitting to see if you need to complete this section. If yes, your manuscript must contain the following sections under the heading `Declarations':

% \begin{itemize}
% \item Funding
% \item Conflict of interest/Competing interests (check journal-specific guidelines for which heading to use)
% \item Ethics approval and consent to participate
% \item Consent for publication
% \item Data availability 
% \item Materials availability
% \item Code availability 
% \item Author contribution
% \end{itemize}

\noindent
% If any of the sections are not relevant to your manuscript, please include the heading and write `Not applicable' for that section. 



\begin{appendices}

\section{Supplementary Tables and Figures}

\begin{table}[h]
\centering
\caption{Supplementary Table S1. Default hyperparameters for all models evaluated in this study.}
\label{tab:supp_s1_hyperparameters}
\begin{tabular}{@{}lll@{}}
\toprule
\textbf{Model} & \textbf{Hyperparameter} & \textbf{Default Value} \\ 
\midrule
\multicolumn{3}{l}{\textit{Tree-based Models}} \\
\midrule
Random Forest & n\_estimators & 100 \\
              & max\_depth & None \\
              & max\_features & sqrt \\
              & min\_samples\_leaf & 1 \\
              & min\_samples\_split & 2 \\
              & bootstrap & True \\
\addlinespace
Quantile Forest & n\_estimators & 300 \\
                & max\_depth & None \\
                & max\_features & sqrt \\
                & min\_samples\_leaf & 1 \\
                & min\_samples\_split & 2 \\
                & bootstrap & True \\
                & quantile & 0.5 \\
\addlinespace
XGBoost & n\_estimators & 100 \\
        & max\_depth & 6 \\
        & learning\_rate & 0.1 \\
        & subsample & 1.0 \\
        & colsample\_bytree & 1.0 \\
        & colsample\_bylevel & 1.0 \\
        & min\_child\_weight & 1 \\
        & gamma & 0.0 \\
        & reg\_alpha & 0.0 \\
        & reg\_lambda & 1.0 \\
\addlinespace
LightGBM & n\_estimators & 100 \\
         & learning\_rate & 0.1 \\
         & num\_leaves & 31 \\
         & max\_depth & $-1$ (unlimited) \\
         & subsample & 1.0 \\
         & colsample\_bytree & 1.0 \\
         & min\_child\_samples & 20 \\
         & reg\_alpha & 0.0 \\
         & reg\_lambda & 0.0 \\
\addlinespace
NGBoost & learning\_rate & 0.01 \\
        & n\_estimators & 500 \\
        & natural\_gradient & True \\
\midrule
\multicolumn{3}{l}{\textit{Kernel-based Models}} \\
\midrule
SVM & C & 1.0 \\
    & gamma & scale \\
    & kernel & rbf \\
\addlinespace
Gauche (GP) & kernel & Tanimoto \\
            & outputscale & 1.0 \\
            & likelihood\_noise & 0.001 \\
\midrule
\multicolumn{3}{l}{\textit{Neural Networks}} \\
\midrule
DNN & hidden\_size\_1 & 128 \\
    & hidden\_size\_2 & 64 \\
    & activation & ReLU \\
    & dropout & 0.2 \\
    & learning\_rate & 0.001 \\
    & batch\_size & 32 \\
\addlinespace
MLP & hidden\_size & 32 \\
    & num\_hidden\_layers & 2 \\
    & dropout\_rate & 0.2 \\
    & learning\_rate & 0.001 \\
\bottomrule
\end{tabular}
\begin{tablenotes}
\small
\item[a] Applied to DNN and MLP architectures.
\end{tablenotes}
\end{table}


To validate the ANOVA design, we computed pairwise Spearman rank correlations between all model NDS profiles and between all representation NDS profiles (Supplementary Tables~\ref{tab:supp_model_redundancy} and~\ref{tab:supp_rep_redundancy}). Models with $\rho > 0.99$ were excluded from the ANOVA to avoid inflating model degrees of freedom with near-duplicate levels. Representations with $\rho > 0.90$ (SNS vs. topological fingerprint) were similarly excluded. We also computed ICC(1,1) for all model pairs to assess within-family consistency (Supplementary Table~\ref{tab:supp_icc}).

\begin{figure}[htbp]
\centering
\includegraphics[width=\textwidth]{fig_methods_noise_strategies_detailed.png}
\caption{Supplementary Figure S1. Detailed noise strategy distributions at multiple $\sigma$ levels showing the progression of label corruption for each strategy.}
\label{fig:supp_noise_detailed}
\end{figure}

\begin{figure}[htbp]
\centering
\includegraphics[width=\textwidth]{fig1_supp_ecfp4_overview.png}
\caption{Supplementary Figure S2. Global overview of noise robustness on the topological fingerprint representation. (A) Performance degradation curves under Gaussian noise. (B) NDS heatmap across model architectures and noise strategies.}
\label{fig:supp_ecfp4_overview}
\end{figure}

\begin{figure}[htbp]
\centering
\includegraphics[width=\textwidth]{fig3_supp_ecfp4_ranking.png}
\caption{Supplementary Figure S3. NDS heatmap and baseline R$^2$ versus NDS scatter for topological fingerprint representation under Gaussian noise.}
\label{fig:supp_ecfp4_ranking}
\end{figure}

\begin{figure}[htbp]
\centering
\includegraphics[width=\textwidth]{fig_full_overview_supp.png}
\caption{Supplementary Figure S4. All model--representation configurations for remaining noise strategies: quantile, heteroscedastic, and value-proportional noise.}
\label{fig:full_overview_supp}
\end{figure}

\begin{figure}[htbp]
\centering
\includegraphics[width=\textwidth]{fig_validation_strategy.png}
\caption{Supplementary Figure S5. Noise strategy effect on robustness across validation datasets. Mean NDS by strategy and dataset.}
\label{fig:supp_validation_strategy}
\end{figure}

% fig_validation_rep_comparison promoted to main text (Figure~\ref{fig:validation_rep_comparison})

\begin{figure}[htbp]
\centering
\includegraphics[width=\textwidth]{fig3_ranking_consistency.png}
\caption{Supplementary Figure S7. Ranking consistency on PDV. (A) NDS heatmap across models and strategies. (B) Baseline R$^2$ versus NDS scatter. (C) Cross-dataset NDS comparison.}
\label{fig:supp_ranking_consistency}
\end{figure}

% fig1_global_overview and fig_interaction promoted to main text (Figures \ref{fig:global_overview} and \ref{fig:interaction})

\begin{figure}[htbp]
\centering
\includegraphics[width=\textwidth]{fig_full_overview.png}
\caption{Supplementary Figure. All model--representation configurations across noise severity tiers. Each point is one model--representation configuration; marker shape indicates representation, color indicates model. Panels show (A) Gaussian, (B) outlier, and (C) threshold noise strategies, spanning the mild-to-severe range. Remaining strategies are shown in Supplementary Figure~\ref{fig:full_overview_supp}.}
\label{fig:supp_full_overview}
\end{figure}

% fig_validation_model_comparison promoted to main text (Figure~\ref{fig:validation_model_comparison})

% Supplementary tables referenced from main text
% All supplementary CSVs are generated by generate_paper_figures.py into results/paper_figures/paper_figures/.
% For Overleaf submission: upload the following CSVs as supplementary data files:
%   - excluded_configs.csv
%   - table1_supp_simple_effects.csv
%   - table5_sigma_rankings.csv
%   - table7_strategy_sensitivity_ratio.csv
%   - table_supp_model_redundancy.csv
%   - table_supp_rep_redundancy.csv
%   - table_supp_icc.csv
%   - table4_supp_uncertainty_by_strategy_rep.csv
%   - table_validation_anova.csv
\begin{table}[h]
\centering
\caption{Supplementary Table. Configurations excluded from robustness analysis due to baseline R$^2 \leq 0.6$. Available as \texttt{excluded\_configs.csv}.}
\label{tab:excluded_configs}
\small
\textit{95 configurations excluded; see supplementary data file.}
\end{table}

\begin{table}[h]
\centering
\caption{Supplementary Table. NDS values across all model--representation--strategy combinations. Available as \texttt{table2\_supp\_nds\_all\_reps.csv}.}
\label{tab:nds_all_reps}
\small
\textit{Full NDS table; see supplementary data file.}
\end{table}

\begin{table}[h]
\centering
\caption{Supplementary Table. Simple effects analysis: $\eta^2$ for model architecture within each representation level, and vice versa, by noise strategy. Available as \texttt{table1\_supp\_simple\_effects.csv}.}
\label{tab:simple_effects}
\small
\textit{Simple effects analysis; see supplementary data file.}
\end{table}

\begin{table}[h]
\centering
\caption{Supplementary Table. Model rankings at each noise level on PDV (Gaussian strategy), showing how relative model performance changes as noise increases. Available as \texttt{table5\_sigma\_rankings.csv}.}
\label{tab:sigma_rankings}
\small
\textit{Sigma-level rankings; see supplementary data file.}
\end{table}

\begin{table}[h]
\centering
\caption{Supplementary Table. Strategy sensitivity ratios relative to Gaussian noise. Each value represents a model's NDS on the given strategy divided by its Gaussian NDS; values $> 1$ indicate greater damage than Gaussian, $< 1$ indicate less.}
\label{tab:supp_sensitivity_ratios}
\small
\begin{tabular}{lccccc}
\toprule
\textbf{Model} & \textbf{Outlier} & \textbf{Quantile} & \textbf{Threshold} & \textbf{Hetero.} & \textbf{Val.-Prop.} \\
\midrule
NGBoost  & 0.33 & 0.87 & 1.97 & 0.55 & 1.75 \\
SVM      & 0.30 & 0.84 & 1.98 & 0.55 & 1.72 \\
LightGBM & 0.32 & 0.86 & 1.98 & 0.55 & 1.75 \\
XGBoost  & 0.35 & 0.88 & 1.98 & 0.54 & 1.77 \\
RF       & 0.32 & 0.86 & 2.00 & 0.54 & 1.76 \\
GP       & 0.34 & 0.87 & 1.91 & 0.56 & 1.71 \\
DNN      & 0.37 & 0.91 & 1.93 & 0.55 & 1.74 \\
MLP      & 0.37 & 0.92 & 1.92 & 0.55 & 1.74 \\
\bottomrule
\end{tabular}
\end{table}

\begin{table}[h]
\centering
\caption{Supplementary Table. Neural network Bayesian transformation effects by representation.}
\label{tab:supp_nn_transforms}
\small
\begin{tabular}{llrrrr}
\toprule
\textbf{Model} & \textbf{Representation} & \textbf{Baseline R$^2$} & \textbf{MLP NDS} & \textbf{BNN-Full NDS} & \textbf{NDS Improvement} \\
\midrule
MLP & Topological & 0.790 & $-0.495$ & $-0.402$ & +0.094 \\
MLP & PDV & 0.828 & $-0.406$ & $-0.369$ & +0.037 \\
MLP & SNS & 0.813 & $-0.546$ & $-0.434$ & +0.112 \\
MLP & SMILES & 0.728 & $-0.600$ & $-0.383$ & +0.217 \\
\bottomrule
\end{tabular}
\end{table}

\begin{table}[h]
\centering
\caption{Supplementary Table. RF vs QRF on validation datasets. QRF is consistently less robust than RF across all external datasets.}
\label{tab:supp_validation_probabilistic}
\small
\begin{tabular}{lrrr}
\toprule
\textbf{Dataset} & \textbf{RF NDS} & \textbf{QRF NDS} & \textbf{$\Delta$ NDS} \\
\midrule
ChEMBL-hERG-Ki     & $-0.279$ & $-0.417$ & $-0.139$ \\
OpenADMET-Caco2    & $-0.555$ & $-0.721$ & $-0.165$ \\
OpenADMET-LogD     & $-0.111$ & $-0.202$ & $-0.091$ \\
\bottomrule
\end{tabular}
\end{table}

\begin{table}[h]
\centering
\caption{Supplementary Table. Pairwise Spearman rank correlations between model NDS profiles. Models with $\rho > 0.99$ were excluded from the ANOVA as near-duplicates. Available as \texttt{table\_supp\_model\_redundancy.csv}.}
\label{tab:supp_model_redundancy}
\small
\textit{Model redundancy analysis; see supplementary data file.}
\end{table}

\begin{table}[h]
\centering
\caption{Supplementary Table. Pairwise Spearman rank correlations between representation NDS profiles. Representations with $\rho > 0.90$ were excluded from the ANOVA. Available as \texttt{table\_supp\_rep\_redundancy.csv}.}
\label{tab:supp_rep_redundancy}
\small
\textit{Representation redundancy analysis; see supplementary data file.}
\end{table}

\begin{table}[h]
\centering
\caption{Supplementary Table. Intraclass correlation coefficients ICC(1,1) for all model pairs, assessing within-family consistency of NDS profiles. Available as \texttt{table\_supp\_icc.csv}.}
\label{tab:supp_icc}
\small
\textit{ICC analysis; see supplementary data file.}
\end{table}

\begin{table}[h]
\centering
\caption{Supplementary Table. Uncertainty quantification metrics broken down by noise strategy and representation for all probabilistic models. Available as \texttt{table4\_supp\_uncertainty\_by\_strategy\_rep.csv}.}
\label{tab:supp_uncertainty_by_strategy_rep}
\small
\textit{Full uncertainty metrics by strategy and representation; see supplementary data file.}
\end{table}

\begin{figure}[htbp]
\centering
\includegraphics[width=\textwidth]{fig_validation_anova.png}
\caption{Supplementary Figure. ANOVA variance decomposition ($\eta^2$, \%) on external validation datasets. Model, representation, and interaction contributions to NDS variance for each dataset.}
\label{fig:supp_validation_anova}
\end{figure}

\begin{table}[h]
\centering
\caption{Supplementary Table. ANOVA variance decomposition on external validation datasets. Available as \texttt{table\_validation\_anova.csv}.}
\label{tab:supp_validation_anova}
\small
\textit{Validation ANOVA; see supplementary data file.}
\end{table}

% An appendix contains supplementary information that is not an essential part of the text itself but which may be helpful in providing a more comprehensive understanding of the research problem or it is information that is too cumbersome to be included in the body of the paper.

%%=============================================%%
%% For submissions to Nature Portfolio Journals %%
%% please use the heading ``Extended Data''.   %%
%%=============================================%%

%%=============================================================%%
%% Sample for another appendix section			       %%
%%=============================================================%%

%% \section{Example of another appendix section}\label{secA2}%
%% Appendices may be used for helpful, supporting or essential material that would otherwise 
%% clutter, break up or be distracting to the text. Appendices can consist of sections, figures, 
%% tables and equations etc.

\end{appendices}
\clearpage

%%===========================================================================================%%
%% If you are submitting to one of the Nature Portfolio journals, using the eJP submission   %%
%% system, please include the references within the manuscript file itself. You may do this  %%
%% by copying the reference list from your .bbl file, paste it into the main manuscript .tex %%
%% file, and delete the associated \verb+\bibliography+ commands.                            %%
%%===========================================================================================%%

\bibliography{sn-bibliography}% common bib file
%% if required, the content of .bbl file can be included here once bbl is generated
%%\input sn-article.bbl


\end{document}
